\section{Discussion}
Across 278 pelvic models, unilateral standing produced larger relative joint motion and bilateral asymmetry than bilateral standing. The localized force-pair responses differed by application site; they did not support one common transverse-expansion pattern. Female translation coefficients were positive in the age-adjusted models, but no residual female translation association survived full adjustment and FDR correction. Negative conditional volume slopes occurred across all five loads. These findings support load-specific structural comparisons, rather than a general female compliance surplus independent of measured anatomy.

\subsection*{Reference configuration and credibility of the endpoints}
Elastic motion was measured against each subject's unloaded mapped anatomy. This distinction matters in population FE models: comparing a shared template directly with a loaded individual would include anatomical variation in the endpoint, potentially inflating associations with geometry. The unloaded-to-unloaded checks and saved-field reconstruction checks provide numerical verification of the extraction, not experimental validation of the equilibrium model.

The affine sacral correction remains a modeling convention rather than only a rigid change of coordinates. Removing sacral strain can change fitted relative motion. The sensitivity analysis found small differences from rigid-only correction in six subjects, but it does not establish cohort-wide external validity. Reported confidence intervals reflect subject sampling and do not cover this extraction choice or FE model discrepancy.

\subsection*{Load dependence and anatomical interpretation}
Unilateral acetabular loading supplies a distinct asymmetric load path, so a paired comparison with bilateral loading directly addresses a mechanical question. However, this static configuration omits the time-varying muscle forces and support conditions of gait. Relative rotation about the ML axis is compatible with sagittal pelvic motion, but its absolute value cannot distinguish nutation from counternutation. Neither ligament tension changes nor frictional self-bracing were measured here. The results therefore do not independently verify force closure, despite its established anatomical relevance \citep{Snijders1993,Vleeming2012}.

The LAB cases apply force pairs at different sites and should be interpreted individually. A reduction in $|d_{CC}|$ accompanied by an increase in $|d_{ML}|$ describes redistribution among the chosen axes. It does not establish increased birth-canal dimensions, preservation of joint congruency or a transition from compression to opening. Such claims would require signed motion in anatomically anchored frames, distances between opposing surfaces and direct inlet, midplane or outlet dimension changes. In particular, ML displacement is not generally normal to the auricular surface. Static force pairs omit fetal contact, pelvic floor mechanics, changing posture and pregnancy-related tissue adaptation \citep{ChenGrimm2021}.

\subsection*{Sex and morphology are associated, not experimentally separated}
The nested regressions show that the initial positive sex coefficients attenuate substantially with measured dimensions; LAB2 changes sign and all three fully adjusted translation intervals include zero. This removes the basis for claiming a persistent female translational surplus independent of the included anatomy. Attenuation is not a percentage of biological mediation: sex, dimensions and unmeasured anatomy are related, covariates are measured with error, and conditional comparisons may have limited support where male and female morphologies overlap poorly. The positive pooled subpubic-angle correlation contrasts with small negative within-sex correlations (Fig.~\ref{fig:sexeffects}), consistent with group separation contributing to the pooled association. No hoop-stiffness mechanism was directly measured. Residual sex coefficients cannot distinguish articular orientation, omitted dimensions, mapping error or other structural differences.

\subsection*{What the material comparison establishes}
Shape-only/full variance ratios were near 100\% for many endpoints but reached 120.61\%, and rotation identity-line $R^2$ fell to 0.878. Thus a reference bone-density field often reproduces the broad between-subject dispersion, but its accuracy is endpoint-dependent and is not uniformly near-perfect. It would not show that mineralization is biologically irrelevant. The implemented modulus law is constant below its density threshold, soft-tissue properties are fixed, and material-only responses are evaluated at one reference geometry. Geometry--material interactions are therefore incompletely sampled. Ligament pretension has been shown to alter SIJ predictions in other FE models \citep{Heyland2025}; its contribution is not bounded by the small bone-density-only variance ratio found here.

Variance ratios above 100\% are mathematically permissible because the variants are separate deterministic systems. They are not percentages of an additive variance budget. Similarly, a high identity-line $R^2$ is an internal agreement measure, not validation against physical measurements. The repaired extraction removes one source of artificial geometry dominance, but external validation and constitutive sensitivity are still required.

\subsection*{Size scaling under prescribed force}
Negative volume slopes need not imply a distinct evolutionary adaptation. For geometrically similar linear-elastic solids with fixed modulus and a fixed resultant force, dimensional scaling gives $u\propto L^{-1}$ and $\theta\sim u/L\propto L^{-2}$, corresponding to volume exponents $-1/3$ and $-2/3$. This is a dimensional benchmark, not a prediction for the present heterogeneous, jointed pelvis with fixed spring parameters. The fitted models additionally condition on linear dimensions, so their volume coefficients describe residual size associations rather than uniform geometric scaling. Forces proportional to body weight or area would give different benchmarks. A nonsignificant sex-by-size interaction does not establish equal slopes; equivalence requires a prespecified tolerance and a suitably powered test.

\subsection*{Originality, usefulness and limits of generalization}
The main potential contribution beyond the earlier morphology study \citep{HenysHammer2025} is a paired map of load-dependent joint responses, coupled to a reproducible comparison of geometry and bone-density variants. Neither a large cohort nor the general importance of geometry is sufficient as a novelty claim. The degree of imaging overlap with previous work must be documented; overlapping anatomy does not constitute independent validation.

The immediate use is methodological: selecting load cases for subsequent experiments, describing distributions of model response, and assessing whether individual bone-density mapping materially changes the selected outputs. No pain data, fixation models, delivery outcomes, parity or pregnancy-specific material measurements were analyzed. The results cannot recommend surgical fixation directions, select a delivery mode, diagnose instability or define normal clinical thresholds.

Absolute magnitudes should be compared with experiments only when loading, reference configuration and endpoint definitions match. Small SIJ movements are challenging to measure, and published estimates depend strongly on methods \citep{Goode2008}. The 16--91-year cohort is not a cohort of laboring women. Its complete records do not guarantee population representativeness. The fixed sacrum, nominal cartilage, linearized ligaments, absence of contact and muscles, RBF approximation, ROI definition and affine correction all constrain interpretation. Mesh convergence, mapped Jacobian quality, stiffness/pretension uncertainty and independent cadaveric or in-vivo validation should accompany any claim of quantitative physiological prediction.
